% Discussion, open problems, acknowledgements, code and data (arXiv conversion: research/paper-review/DECISIONS.md
% §A4, §A10, §D1, §F18, §F23; framing: REFRAME.md §4). The open problems list refers to every question and conjecture of the paper.
% Rename of 2026-10-15 (research/paper-review/rename/RENAME.md §4, Paper 1 item 7): names updated, versions named where they
% matter; the open problems are unchanged. Log: research/paper-review/rename/log-p1-front.md.
\section{Discussion}\label{sec:discussion}

\paragraph{SOI as ideal supervised learning.} \Cref{rem:learn:ideal} states in what sense SOI is ideal supervised learning and where the analogy stops: its guarantees compare it with computable labellers, in log loss, and need no assumption on how the inputs are chosen. The out-of-distribution results have the same character (\cref{rem:ood:lopsided,rem:ood:answer}): on inputs that a simple property separates from the training data, SOI keeps a probability bounded below, depending on the universal machine, for every simple behaviour, whatever the amount of data (\cref{prop:ood:hanni,prop:ood:machine}); and for a computable target $f$ its total posterior error is at most $(\ln2)(K(f)+c)$ (\cref{thm:ood:cumulative}).

\paragraph{What the hypotheses may see.} The large gaps between SOI and Solomonoff induction in \cref{sec:hist,sec:ex} are properties of order-free prediction rather than of any one version. Every order-independent predictor, Solomonoff's per-query solution included, pays them in expectation (\cref{prop:hist:orderfree,prop:hist:operator,prop:ex:orderfree,prop:ex:law}), and hypotheses that see the list of previous inputs remove them (\cref{cor:hist:H4,prop:ex:cychist,thm:ex:eisensi}). Which of the order-free versions and \SFIhist{} is the right ideal therefore depends on whether the order of the data carries information about the labels. With i.i.d.\ inputs it does not, and the guarantees of SOI apply as stated. When a teacher chooses the order, it can carry almost all of the information, as in the cyclic permutation, where the labels repeat the inputs. The set of previous inputs is not enough to recover this information (\cref{prop:hist:strict,prop:ex:cycset}), and even with the list, the loss of \SFIhist{} is at most SI's joint loss plus a constant, not automatically at most SI's conditional label loss plus a constant (\cref{thm:hist:si,prop:hist:inputloss}).

\paragraph{Which version to take as the default.} The versions of \cref{sec:versions} agree up to $K(\dom D)+O(1)$ bits on data sets (Solomonoff's per-query solution on distinct inputs), and their gaps are unbounded. Between the program-based versions and the one-tape version they are small on the known witnesses: about $\log d$ for the Kolmogorov version, the single shortest program, on two points with $d$-bit inputs, and a padded tape-length cost for the coin-free version, the program mixture (\cref{cor:versions:twopoint,prop:versions:tape}); the weak semimeasure, which has no posterior over functions, can exceed $\MF$ by linearly many bits (\cref{thm:setting:weakvsmachine}). For deterministic computable targets all of them learn at the same rate (\cref{rem:versions:prediction,prop:learn:occam}), so the choice matters mostly for random labels and for code lengths. We take the one-tape version $\MF$ as the default because it treats the labelling as one random function and is the law of one universal machine with one sequence of coin flips, it has a posterior over functions, it dominates every machine law with constant $2^{-K(\nu)}$, it is invariant, and its semimeasure-convention loss does not depend on the order of the data (\cref{sec:setting:sfi,rem:setting:recommend}). Where the randomness lives also matters for noise: a random function, such as the one-tape version, cannot produce fresh noise on a repeated input, while Solomonoff's per-query solution can (\cref{prop:learn:freshnoise,prop:learn:perquery}).

\paragraph{Computable versions and networks.} Time-bounded mixtures keep constant regret against all predictors that see only the input and run within a time bound (\cref{prop:comp:round,cor:comp:poly}). For circuits, Bayes mistakes and circuit size are polynomially related on every ordering (\cref{thm:comp:circuitbayes}), while simplest-consistent learning on the same prior makes $2^{\Omega(\ell/\log\ell)}$ mistakes, exponentially many in the input dimension, on some ordering for the conjunction of $d$ variables, whose description length is $\ell=\Theta(d\log d)$ (\cref{thm:comp:ladder}); with the same prior the posterior-majority predictor makes at most $\ell$ mistakes on every ordering, so it is the selection rule, not the prior, that fails (\cref{rem:comp:majority,rem:comp:selection}). For neural networks, the parameter--function map alone implies no simplicity bias (\cref{prop:nn:nolower}); any bridge to circuit SOI must come from the architecture and the prior, which is what the slides of \citet{vaintrobhanni2026nnbayes} claim for small prior scales (\cref{rem:nn:slides}). In the Gaussian-process limits, kernel regression and Gaussian-process posterior means need about $\binom dk$ samples, in squared error, on random $k$-sparse parities where SOI needs $O(k\log d)$ for any fixed accuracy, and bias-free two-layer ReLU kernels never learn odd parities of degree at least $3$ (\cref{cor:nn:separation,rem:nn:oddparity}); whether NN Bayes generalises like SGD-trained networks is not addressed here.

\paragraph{The companion paper.} Where the randomness lives also matters for steeper simplicity priors (\cref{sec:steep}). The companion paper studies Solomonoff's per-query solution under such priors and gives two reasons for them \citep{claudeSFI2026steeper}. With private randomness, a penalty steeper than likelihood by a fixed factor splits per-query samples into the sampler and its noise, exactly in the limit and at a rate that depends only on the factor, the program's length, a margin and the confidence. A penalty steeper by a factor proportional to the sample size, in an explicit window, keeps hypotheses that are near-optimal for the population objective from modelling a teacher's errors on a parity task when the errors are made at a rate below $0.05596$ on an error set of product form that is incompressible relative to the parity. With shared randomness, as in the one-tape version, or in sequence prediction, a universal mixture is itself a hypothesis and every selected hypothesis is shorter than a bound that does not depend on the data. Whether structured error rules leave a window for steeper selection is open there.

\subsection{Open problems}\label{sec:discussion:open}

\begin{enumerate}[leftmargin=*,label=(\arabic*)]
\item Is every lower-semicomputable law of a random partial function dominated by $\MF$ up to a constant factor (\cref{q:setting:rpf})? Is there a two-sided characterisation of $\LSFI(D)$ by standard complexities (\cref{q:setting:twosided})? What is the arithmetical level of the conditionals of SOI (\cref{q:setting:level})?
\item Is the normalised loss of $\MF$ order-dependent, and which lower bounds proved here only in the semimeasure convention hold for it (\cref{q:setting:normorder}), in particular over the full cyclic permutation and in Eisenstat's example (\cref{q:ex:normalised})?
\item Is $\log(\MF(D)/W(D))=O(K(d,|D|))$ for inputs of length at most $d$, and what are the exact rates of the gaps of \cref{tab:versions:compare} (\cref{q:versions:open})? Does some gap tend to infinity along a computable family of data sets (\cref{q:versions:explicit})?
\item Does the last iterate, not only the average, have expected risk $O(K(f)/n)$ (\cref{rem:learn:lastiterate})? Is the leave-one-out error of the one-tape version $O(\KF(f|_B)/|B|)$ up to logarithmic factors (\cref{conj:learn:loo})?
\item For computable reveal laws, can the term $K(n)+K(A)$ of \cref{cor:hist:thm1} be dropped, and do defective samplers, infinite alphabets and reveal laws that see $f$ admit the same bound (\cref{rem:hist:thm1open,rem:hist:minimaxsketch})? Can $K(\Delta_f)$ be replaced by the algorithmic mutual information of $\Delta_f$ and $f$ (\cref{conj:hist:algmi})? How does \SFIhist{} compare with normalised Solomonoff induction, and do the separations of \cref{prop:hist:strict,thm:hist:gf} hold for normalised SOI (\cref{q:hist:open})?
\item Does the depth count for the cyclic permutation hold for every weak function semimeasure (\cref{q:ex:depthweak})? With raw, undelimited blocks, is SI's expected label loss on the cyclic permutation at least $K(d)-\Kpad(d)-O(1)$ (\cref{conj:ex:raw})?
\item Is there a positive out-of-distribution statement for properties that are simple given a program fitted to the training data (\cref{q:ood:converse})? Is there a natural universal machine for which the intended answers of \cref{ex:ood:hanniexamples} get posterior probability close to $1$ (\cref{conj:ood:natural})?
\item Is there a learner with time per round polynomial in the time bound and constant regret against all time-bounded samplers (\cref{q:comp:samplers})? Is there a circuit coding theorem for size priors with polynomial tails, and are there Occam failures with $2^{\Omega(\ell)}$ mistakes (\cref{q:comp:open})?
\item Is the uniform prior over circuits of a given size bounded above by the circuit simplicity prior up to polynomial factors, is there a lower bound on the log prior under NN Bayes at a fixed prior scale, and does the main theorem of the slides hold as stated (\cref{q:nn:bridges})?
\end{enumerate}

\section*{Acknowledgements}\addcontentsline{toc}{section}{Acknowledgements}
The coin example of \cref{sec:ex:eisenstat} is due to Sam Eisenstat (personal communication). \Cref{sec:nn} and the conjunction ladder of \cref{thm:comp:ladder} build on joint work with Dmitry Vaintrob \citep{vaintrobhanni2026nnbayes}.

Several of the questions and examples studied here first appeared in H\"anni's working notes \citep{hanni2026notes}. Further, independent instances of Claude checked each part: they re-derived the proofs and re-implemented the computations, and the results were revised in response. No human has yet independently checked all of the mathematics, and readers should weigh the results accordingly. H\"anni takes responsibility for the content of this paper.

\section*{Code and data}\addcontentsline{toc}{section}{Code and data}
All computations reported here were done with seeded or deterministic Python scripts (Python~3.11 with NumPy, SciPy and mpmath). The scripts, their saved outputs and independent re-implementations of the main computations are in the folder \texttt{solomonoff-operator-induction/} of the repository \url{https://github.com/kaarelh/AI-works}; its README lists the script behind each reported number and how to re-run them. Re-running the scripts reproduces the saved outputs, apart from timing information and one time-limited optimisation.
